\documentclass[../../main.tex]{subfiles}
\begin{document}

\begin{center}
  {\fontsize{48pt}{58pt}\selectfont \textbf{2009}\normalsize\textbf{年度}}
\end{center}

\begin{table}[htb]
  \centering
  \caption{出題テーマと難易度}
  \begin{tabular}{|c|p{9cm}|c|}
    \hline
    \textbf{問題番号} & \textbf{テーマ} & \textbf{難易度} \\
    \hline
    第1問 & 【乙】直方体と垂線の交点／【甲】直方体の座標 & \\
    \hline
    第2問 & 【乙】三角形と円の中心／【甲】対称な三角形 & \\
    \hline
    第3問 & 【乙】カードの移動確率／【甲】対数不等式 & \\
    \hline
    第4問 & 行列の反復作用とベクトルの長さ & \\
    \hline
    第5問 & 【乙】極座標曲線の回転体／【甲】階乗の素因数 & \\
    \hline
    第6問 & 【乙】二次体の冪と互いに素な係数／【甲】極座標曲線の長さ & \\
    \hline
  \end{tabular}
\end{table}

\begin{table}[htb]
  \centering
  \caption{解答}
  \renewcommand{\arraystretch}{1.4}
  \begin{tabular}{|c|c|p{9cm}|}
    \hline
    \textbf{問題番号} & \textbf{小問} & \textbf{答え} \\
    \hline
    第1問 & 【乙】 & $\dfrac{16}{9}s+4t=1$，$0<s<9/16$，$0<t<1/4$ \\
    \cline{2-3}
    & 【甲】 &  \\
    \hline
    第2問 & 【乙】 &  \\
    \cline{2-3}
    & 【甲】 &  \\
    \hline
    第3問 & 【乙】 & $\dfrac{(n-1)!(n^2-n+2)}{2n^n}$ \\
    \cline{2-3}
    & 【甲】 &  \\
    \hline
    第4問 & -- &  \\
    \hline
    第5問 & 【乙】 & $\dfrac{40\pi}{3}$ \\
    \cline{2-3}
    & 【甲】 &  \\
    \hline
    第6問 & 【乙】(1) &  \\
    \cline{2-3}
    & 【乙】(2) &  \\
    \cline{2-3}
    & 【甲】 &  \\
    \hline
  \end{tabular}
\end{table}

\paragraph{試験形態}
6問．試験時間は要確認．

\end{document}
